\begin{theorem}[Affine blow-ups at a fixed center]
Let $E$ be a finite-dimensional real inner-product space, let $D\subseteq E$,
let $w:E\to\R$ be infinity-harmonic on $D$ in the sense of
\noderef{IsInfHarmonicOn}, and let $z\in D$. For a sequence of radii
$\rho_i>0$ with $\rho_i\to0$ put
\[
  V_i(Y)=\frac{w(z+\rho_iY)-w(z)}{\rho_i}\qquad(Y\in E).
\]
Then:
\begin{enumerate}
\item every $V:E\to\R$ which is the limit of such a family $V_i$, uniformly on
every compact subset of $E$, is a continuous linear function on $E$; and
\item for every sequence of radii $\rho_i>0$ with $\rho_i\to0$ there are a
strictly increasing sequence of indices $i\mapsto\iota(i)$ and a function
$V:E\to\R$ such that $V_{\iota(i)}\to V$ uniformly on every compact subset of
$E$.
\end{enumerate}
\end{theorem}
\begin{proof}
By \noderef{IsInfHarmonicOn} the set $D$ is open and $w$ is continuous on $D$,
so there is $\sigma>0$ with $\overline B_{2\sigma}(z)\subseteq D$.

\emph{Step 1: $w$ is Lipschitz on $\overline B_\sigma(z)$.}
The set $\overline B_{2\sigma}(z)$ is compact and contained in $D$, so
$\operatorname{osc}=\max_{\overline B_{2\sigma}(z)}w
 -\min_{\overline B_{2\sigma}(z)}w$ is finite; put
$\Lambda=\operatorname{osc}/\sigma\ge0$. Then
$|w(a)-w(b)|\le\operatorname{osc}=\Lambda\sigma$ for
$a,b\in\overline B_{2\sigma}(z)$, so \noderef{InfHarmonicLocLipschitz}, applied
with $z_0=z$, $\rho=\sigma$, the same $\sigma$ and this $\Lambda$ — its
hypotheses $\overline B_{\rho+\sigma}(z)=\overline B_{2\sigma}(z)\subseteq D$ and
the oscillation bound hold — gives
\begin{equation}\label{eq:ab-lip}
  |w(a)-w(b)|\le\Lambda|a-b|\qquad\bigl(a,b\in\overline B_\sigma(z)\bigr).
\end{equation}

\emph{Step 2: uniform bounds for the blow-ups.}
Let $\rho_i>0$ with $\rho_i\to0$ and put
$V_i(Y)=\bigl(w(z+\rho_iY)-w(z)\bigr)/\rho_i$. Fix $R>0$. For $|Y|\le R$ and all
$i$ with $\rho_iR\le\sigma$ we have $z+\rho_iY\in\overline B_\sigma(z)$, so by
\eqref{eq:ab-lip} with $a=z+\rho_iY$, $b=z$,
\begin{equation}\label{eq:ab-unif}
  |V_i(Y)|\le\Lambda|Y|\le\Lambda R,
  \qquad
  |V_i(Y)-V_i(Y')|\le\Lambda|Y-Y'|\qquad(|Y|,|Y'|\le R),
\end{equation}
the second bound by \eqref{eq:ab-lip} with $a=z+\rho_iY$, $b=z+\rho_iY'$. Since
$\rho_i\to0$, all but finitely many $i$ satisfy $\rho_iR\le\sigma$. Note also
$V_i(0)=0$ for every $i$.

\emph{Step 3: proof of (2).}
The selection of a locally uniformly convergent subsequence needs a compactness
statement which is not among the background items of the preamble, so we prove
the instance we use; the argument is the standard Arzel\`a--Ascoli one, and the
only external inputs are the preamble items \emph{Compactness, Weierstrass and
uniform continuity} (a closed bounded subset of a finite-dimensional space is
compact) and \emph{Bolzano--Weierstrass}.

\emph{Claim A.} Let $\mathcal K\subseteq E$ be closed and bounded, let $A\ge0$
and $\Lambda\ge0$, let $S\subseteq\mathbb N$ be infinite and let $g_i:E\to\R$,
$i\in S$, satisfy
\[
  |g_i(a)|\le A,\qquad|g_i(a)-g_i(b)|\le\Lambda|a-b|
  \qquad(a,b\in\mathcal K,\ i\in S).
\]
Then there are an infinite set $S'\subseteq S$ and a function $g$ on
$\mathcal K$ such that $\sup_{\mathcal K}|g_i-g|\to0$ as $i\to\infty$ along
$S'$.

\emph{Proof of Claim A.} If $\mathcal K=\emptyset$ take $S'=S$ and $g$ the empty
function; so assume $\mathcal K\ne\emptyset$. Being closed and bounded in the
finite-dimensional space $E$, the set $\mathcal K$ is compact. For each integer
$k\ge1$ the balls $B_{1/k}(a)$, $a\in\mathcal K$, form an open cover of
$\mathcal K$, so finitely many of them cover it; let $N_k\subseteq\mathcal K$ be
a finite nonempty set with
$\mathcal K\subseteq\bigcup_{b\in N_k}B_{1/k}(b)$, and put
$Q=\bigcup_{k\ge1}N_k$. Then $Q\subseteq\mathcal K$ is countable and nonempty,
and for every $\varsigma>0$ the balls $B_\varsigma(b)$, $b\in Q$, cover
$\mathcal K$: choose $k$ with $1/k\le\varsigma$ and use $N_k$. Enumerate
$Q=\{a_1,a_2,\dots\}$ (if $Q$ is finite the argument below is the same and
shorter).

For each $k$ the family $\bigl(g_i(a_k)\bigr)_{i\in S}$ is bounded by $A$, so by
Bolzano--Weierstrass we may choose nested infinite sets
$S\supseteq J_1\supseteq J_2\supseteq\cdots$ such that for every $k$ the values
$g_i(a_k)$ converge as $i\to\infty$ along $J_k$: given $J_{k-1}$ (with
$J_0=S$), pick a convergent subsequence of $\bigl(g_i(a_k)\bigr)_{i\in J_{k-1}}$
and let $J_k$ be its set of indices. Choose $i_1<i_2<\cdots$ with $i_l\in J_l$,
which is possible because each $J_l$ is infinite, and put
$S'=\{i_l:l\ge1\}$, an infinite subset of $S$. Since
$\{i_l:l\ge k\}\subseteq J_k$ for every $k$, the sequence
$\bigl(g_{i_l}(a_k)\bigr)_l$ converges for every $k$.

Fix $\varepsilon>0$ and put $\eta=\varepsilon/\bigl(3(\Lambda+1)\bigr)>0$. Cover
$\mathcal K$ by finitely many balls $B_\eta(b_1),\dots,B_\eta(b_N)$ with
$b_\nu\in Q$, as in the first paragraph, and choose $l_0$ with
$|g_{i_l}(b_\nu)-g_{i_{l'}}(b_\nu)|\le\varepsilon/3$ for all $\nu\le N$ and all
$l,l'\ge l_0$; this is possible because each of the finitely many sequences
$\bigl(g_{i_l}(b_\nu)\bigr)_l$ converges, hence is Cauchy. Given
$a\in\mathcal K$, pick $\nu$ with $|a-b_\nu|<\eta$; then for $l,l'\ge l_0$,
\[
  |g_{i_l}(a)-g_{i_{l'}}(a)|
   \le|g_{i_l}(a)-g_{i_l}(b_\nu)|
     +|g_{i_l}(b_\nu)-g_{i_{l'}}(b_\nu)|
     +|g_{i_{l'}}(b_\nu)-g_{i_{l'}}(a)|
   \le2\Lambda\eta+\frac\varepsilon3\le\varepsilon ,
\]
using $2\Lambda\eta\le2\varepsilon/3$. So $\bigl(g_{i_l}\bigr)_l$ is uniformly
Cauchy on $\mathcal K$; for each $a\in\mathcal K$ the real sequence
$\bigl(g_{i_l}(a)\bigr)_l$ is therefore Cauchy and has a limit $g(a)$, and
letting $l'\to\infty$ in the display gives
$\sup_{\mathcal K}|g_{i_l}-g|\le\varepsilon$ for every $l\ge l_0$. As
$\varepsilon>0$ was arbitrary, $\sup_{\mathcal K}|g_i-g|\to0$ along $S'$. This
proves Claim A.

We now prove (2). Let $\rho_i>0$ with $\rho_i\to0$. For each integer $R\ge1$
there is, since $\rho_i\to0$, an index $m_R$ with $\rho_iR\le\sigma$ for all
$i\ge m_R$; for such $i$ the bounds \eqref{eq:ab-unif} give
$|V_i|\le\Lambda R$ and $|V_i(Y)-V_i(Y')|\le\Lambda|Y-Y'|$ on
$\overline B_R(0)$, a closed bounded set. Define infinite index sets
$S_1\supseteq S_2\supseteq\cdots$ recursively: $S_R$ is the set $S'$ produced by
Claim A applied with $\mathcal K=\overline B_R(0)$, $A=\Lambda R$, this
$\Lambda$, and the index set $\{i\ge m_1\}$ when $R=1$ and
$\{i\in S_{R-1}:i\ge m_R\}$ when $R\ge2$; these index sets are infinite and
their elements $i$ all satisfy $\rho_iR\le\sigma$, so the hypotheses of Claim A
hold. Claim A also supplies functions $V^{(R)}$ on $\overline B_R(0)$ with
$\sup_{\overline B_R(0)}|V_i-V^{(R)}|\to0$ as $i\to\infty$ along $S_R$.

Choose $\iota(1)<\iota(2)<\cdots$ with $\iota(j)\in S_j$ for every $j$, which is
possible because each $S_j$ is infinite; thus $\iota:\mathbb N\to\mathbb N$ is
strictly increasing. Fix an integer $R\ge1$. For $j\ge R$ we have
$\iota(j)\in S_j\subseteq S_R$, and $\iota(j)\to\infty$, so
\begin{equation}\label{eq:ab-diag}
  \sup_{\overline B_R(0)}\bigl|V_{\iota(j)}-V^{(R)}\bigr|\longrightarrow0
  \qquad(j\to\infty).
\end{equation}
In particular, for $R\le R'$ the functions $V^{(R)}$ and $V^{(R')}$ agree on
$\overline B_R(0)$, both being the pointwise limit of $V_{\iota(j)}$ there.
Hence a single function $V$ on $E$ is well defined by $V(Y)=V^{(R)}(Y)$ for any
integer $R\ge\max\{1,|Y|\}$, and \eqref{eq:ab-diag} says that
$V_{\iota(j)}\to V$ uniformly on $\overline B_R(0)$ for every integer $R\ge1$.
Since every compact $K\subseteq E$ is bounded, hence contained in some
$\overline B_R(0)$ with $R$ an integer, $R\ge1$, the convergence is uniform on
every compact subset of $E$. This is (2).

\emph{Step 4: $V$ is infinity-harmonic on $E$ and $\Lambda$-Lipschitz.}
From now on let $\rho_i>0$, $\rho_i\to0$, and let $V$ be a function on $E$ such
that $V_i\to V$ uniformly on every compact subset of $E$, as in the hypothesis
of (1).

By \noderef{InfHarmonicRescale}, applied with $x_0=z$ and $A=\rho_i>0$, the
function $Y\mapsto w(z+\rho_iY)$ is infinity-harmonic on
$D_i=\{Y:z+\rho_iY\in D\}$; and by \noderef{InfHarmonicNormalize}, applied with
$c=w(z)$ and $M=\rho_i>0$ on the open set $D_i$, the function $V_i$ is
infinity-harmonic on $D_i$. Moreover $\overline B_\sigma(z)\subseteq D$ gives
$B_{\sigma/\rho_i}(0)\subseteq D_i$.

Fix $R>0$ and choose $i_0$ with $\sigma/\rho_i>R$ for $i\ge i_0$; then each
$V_{i_0+j}$, $j\ge0$, is infinity-harmonic on the open ball $B_R(0)$. Indeed
$B_R(0)\subseteq D_i$ is open, and the restriction of the infinity-harmonicity
of \noderef{IsInfHarmonicOn} to an open subset is legitimate: continuity
restricts, and if $\Omega''\subseteq\Omega'$ is open, $\zeta\in\Omega''$ and
$V_i-\phi$ has a local maximum on $\Omega''$ at $\zeta$, then it has one on
$\Omega'$ at $\zeta$, because a neighbourhood $N$ of $\zeta$ with
$(V_i-\phi)\le(V_i-\phi)(\zeta)$ on $N\cap\Omega''$ may be replaced by
$N\cap\Omega''$, itself a neighbourhood of $\zeta$, on which
$N\cap\Omega''\cap\Omega'=N\cap\Omega''$; the same applies to minima.

The shifted sequence $j\mapsto V_{i_0+j}$ still converges to $V$ uniformly on
every compact subset of $B_R(0)$, being a tail of the original one. Hence
\noderef{InfHarmonicStability}, applied with $\Omega=B_R(0)$, shows that $V$ is
infinity-harmonic on $B_R(0)$. As $R>0$ was arbitrary and a local extremum of
$V-\phi$ on $E$ at a point $\zeta$ is in particular one on $B_R(0)$ for any
$R>|\zeta|$, together with the continuity of $V$ on each $B_R(0)$, we conclude
that $V$ is infinity-harmonic on $E$.

Finally, letting $i\to\infty$ in \eqref{eq:ab-unif} gives
\begin{equation}\label{eq:ab-Vlip}
  |V(Y)-V(Y')|\le\Lambda|Y-Y'|\quad(Y,Y'\in E),\qquad V(0)=0 .
\end{equation}

\emph{Step 5: the slopes of $V$ at the origin.}
By \noderef{ConeSlopeLimitEq}, applied to $w$ on $D$ at $z\in D$, there is
$L_*\ge0$ with
\begin{equation}\label{eq:ab-Lstar}
  \lim_{s\downarrow0}L_s^+(w;z)=\lim_{s\downarrow0}L_s^-(w;z)=L_* .
\end{equation}
Fix $R>0$. By \noderef{ConeSlopeUpper}, for $\rho_iR\le\sigma$,
\[
  \sup_{|Y|\le R}\frac{V_i(Y)-V_i(0)}{R}
   =\frac{1}{R\rho_i}\sup_{|h|\le\rho_iR}\bigl(w(z+h)-w(z)\bigr)
   =L_{\rho_iR}^+(w;z),
\]
where we substituted $h=\rho_iY$, which is a bijection of $\{|Y|\le R\}$ onto
$\{|h|\le\rho_iR\}$, and used $V_i(0)=0$. The left side is
$L_R^+(V_i;0)$ by \noderef{ConeSlopeUpper}. Since $V_i\to V$ uniformly on
$\overline B_R(0)$ and all these suprema are finite by \eqref{eq:ab-unif}, we
get $L_R^+(V_i;0)\to L_R^+(V;0)$; and $\rho_iR\to0$, so by
\eqref{eq:ab-Lstar} the right side tends to $L_*$. Hence, and by the identical
computation with $-(w(z+h)-w(z))$ and \noderef{ConeSlopeLower},
\begin{equation}\label{eq:ab-origin}
  L_R^+(V;0)=L_R^-(V;0)=L_*\qquad(R>0).
\end{equation}
In particular, for $Y\ne0$ the competitor $Y$ at radius $R=|Y|$ in
\noderef{ConeSlopeUpper} and \noderef{ConeSlopeLower} gives
$V(Y)/|Y|\le L_*$ and $-V(Y)/|Y|\le L_*$, so
\begin{equation}\label{eq:ab-cone}
  |V(Y)|\le L_*|Y|\qquad(Y\in E),
\end{equation}
the case $Y=0$ being $V(0)=0$.

\emph{Step 6: $V$ is globally $L_*$-Lipschitz.}
Let $q\in E$ and $T>0$. By \noderef{ConeSlopeUpper} and \eqref{eq:ab-cone},
\[
  L_T^+(V;q)=\frac1T\sup_{|H|\le T}\bigl(V(q+H)-V(q)\bigr)
   \le\frac{L_*(T+|q|)+|V(q)|}{T},
\]
and the same bound holds for $L_T^-(V;q)$ by \noderef{ConeSlopeLower}. Fix
$T_0>0$. Since $V$ is infinity-harmonic on $E$ by Step 4 and
$\overline B_T(q)\subseteq E$ for every $T$, the radial monotonicity assertion of \noderef{ConeComparison}, applied at the
center $q$ with radius $T_0$ and comparison radius $T\ge T_0$, gives
$L_{T_0}^\pm(V;q)\le L_T^\pm(V;q)$. Letting $T\to\infty$ in the displayed bound
yields
\begin{equation}\label{eq:ab-allslopes}
  L_{T_0}^+(V;q)\le L_*,\qquad L_{T_0}^-(V;q)\le L_*
  \qquad(q\in E,\ T_0>0).
\end{equation}
Now let $Q,q\in E$ with $Q\ne q$ and put $T_0=|Q-q|>0$. The two cone comparison
inequalities of \noderef{ConeComparison} at the center $q$ with radius $T_0$
give
$V(Q)-V(q)\le L_{T_0}^+(V;q)|Q-q|$ and
$-(V(Q)-V(q))\le L_{T_0}^-(V;q)|Q-q|$, so by \eqref{eq:ab-allslopes}
\begin{equation}\label{eq:ab-globlip}
  |V(Q)-V(q)|\le L_*|Q-q|\qquad(Q,q\in E).
\end{equation}

\emph{Step 7: the case $L_*=0$.}
If $L_*=0$ then \eqref{eq:ab-cone} gives $V\equiv0$, which is the zero
continuous linear functional, and (1) holds.

\emph{Step 8: the case $L_*>0$: a distinguished direction.}
Assume $L_*>0$. For $R>0$ the supremum in $L_R^+(V;0)$ is over the compact ball
$\overline B_R(0)$ of the continuous function $Y\mapsto V(Y)/R$, so it is
attained: by \eqref{eq:ab-origin} there is $x_R$ with $|x_R|\le R$ and
$V(x_R)=L_*R$. Likewise there is $y_S$ with $|y_S|\le S$ and
$-V(y_S)=L_*S$.

For any $R,S>0$, using \eqref{eq:ab-globlip},
\begin{equation}\label{eq:ab-equality}
  L_*(R+S)=V(x_R)-V(y_S)\le L_*|x_R-y_S|\le L_*\bigl(|x_R|+|y_S|\bigr)
   \le L_*(R+S).
\end{equation}
Since $L_*>0$, every inequality in \eqref{eq:ab-equality} is an equality;
hence $|x_R|=R$, $|y_S|=S$ and $|x_R-y_S|=|x_R|+|y_S|$.

Equality in $|a-b|\le|a|+|b|$ for nonzero $a,b$ forces $a$ and $-b$ to be
positive multiples of a common unit vector. Indeed, squaring
$|a-b|^2=(|a|+|b|)^2$ gives $-2\langle a,b\rangle=2|a||b|$, i.e.
$\langle a,-b\rangle=|a||{-b}|$, which is equality in Cauchy--Schwarz and
therefore means $-b=(|b|/|a|)\,a$.

Fix $S=1$ and let $e=-y_1$, a unit vector since $|y_1|=1$. For each $R>0$,
applying the previous paragraph with $a=x_R$ and $b=y_1$ gives
$e=-y_1=(1/R)x_R$, i.e. $x_R=Re$, so
\begin{equation}\label{eq:ab-ray-pos}
  V(Re)=L_*R\qquad(R>0).
\end{equation}
Now fix $R=1$; the same paragraph with $a=x_1=e$ and $b=y_S$ gives
$-y_S=S\,e$, i.e. $y_S=-Se$, so
\begin{equation}\label{eq:ab-ray-neg}
  V(-Se)=-L_*S\qquad(S>0).
\end{equation}
Together with $V(0)=0$, \eqref{eq:ab-ray-pos} and \eqref{eq:ab-ray-neg} give
\begin{equation}\label{eq:ab-ray}
  V(\tau e)=L_*\tau\qquad(\tau\in\R).
\end{equation}

\emph{Step 9: $V$ is linear.}
Let $Y\in E$. By \eqref{eq:ab-globlip} and \eqref{eq:ab-ray}, for every
$\tau\in\R$,
\[
  L_*\tau-L_*|Y-\tau e|\le V(Y)\le L_*\tau+L_*|Y-\tau e| .
\]
For $\tau>0$ we compute
$|Y-\tau e|^2=\tau^2-2\tau\langle e,Y\rangle+|Y|^2$, so
\[
  |Y-\tau e|=\tau\sqrt{1-\frac{2\langle e,Y\rangle}{\tau}+\frac{|Y|^2}{\tau^2}}
   =\tau-\langle e,Y\rangle+O\bigl(\tau^{-1}\bigr)
   \qquad(\tau\to+\infty),
\]
by the expansion $\sqrt{1+u}=1+u/2+O(u^2)$ applied with
$u=-2\langle e,Y\rangle/\tau+|Y|^2/\tau^2\to0$. Hence
$L_*\tau-L_*|Y-\tau e|\to L_*\langle e,Y\rangle$ as $\tau\to+\infty$, and the
lower bound gives $V(Y)\ge L_*\langle e,Y\rangle$. Taking instead
$\tau=-\varsigma$ with $\varsigma\to+\infty$, the same expansion with $e$
replaced by $-e$ gives
$|Y+\varsigma e|=\varsigma+\langle e,Y\rangle+O(\varsigma^{-1})$, so
\[
  L_*\tau+L_*|Y-\tau e|=-L_*\varsigma+L_*|Y+\varsigma e|
   \longrightarrow L_*\langle e,Y\rangle ,
\]
and the upper bound gives $V(Y)\le L_*\langle e,Y\rangle$. Therefore
\[
  V(Y)=L_*\langle e,Y\rangle\qquad(Y\in E),
\]
so $V$ is the continuous linear functional $L=L_*\langle e,\cdot\rangle$ on $E$.
This proves (1).
\end{proof}
